Production--logistics collaborative scheduling couples a flexible job shop
with a fleet of automated guided vehicles. We formulate the two layers as a
single decision chain and train it with group-relative policy optimization, so
that every production and logistics decision in an episode contributes to one
log-probability and one advantage. The policy is a standard Transformer over a
token sequence with a candidate-scoring head per decision kind, and ten
constraints are modelled as independent switches, each with a mechanism that
gives the policy something to act on. We measure the method on MK01 of the MKT
benchmark against a rule baseline, a single-objective version of itself, and
NSGA-II. On the fully constrained problem it dominates all twelve
out-of-sample points of the metaheuristic front; on the unconstrained problem
the metaheuristic is marginally ahead. Three further results are reported.
First, the aggregate energy consumption commonly reported in this literature is
close to rank-equivalent with the makespan, because the terms that dominate it
are linear in the horizon by construction, so energy adds a third dimension
only under a net account that removes those terms. Second, an ablation over the
advantage shows that subtracting a group-mean baseline has a measurable effect
while dividing by the group standard deviation does not. Third, a hypothesis
that single-objective ablation systematically misses mechanisms did not survive
its test.
